\section{Introduction}
Reinforcement learning (RL) has shown exceptional potential in solving sequential decision-making problems.
However, in real-world scenarios involving complex, long-term tasks with sparse rewards, agents face significant challenges in planning and decision-making due to insufficient feedback signals\cite{sutton1998reinforcement, dulac2021challenges}.
Traditional RL frameworks address these challenges primarily by focusing on two strategies: task decomposition and reward shaping.
Hierarchical RL\cite{nachum2018data, pateria2021hierarchical} simplifies decision-making by autonomously decomposing complex tasks into hierarchical subtasks, while reward shaping techniques provide dense training signals through the reconstruction of reward functions, such as inverse RL\cite{arora2021survey, ziebart2008maximum, finn2016guided} and intrinsic reward mechanisms\cite{wan2023deir}.
Despite their advantages, both approaches heavily rely on expert data for priors\cite{penzkofer2024int} or supervision\cite{pan2018human}, and the large-scale acquisition of such data is often challenging, limiting the scalability and applicability of these methods.

Recent advancements in large language model (LLM) research have made significant progress in areas such as dialogue\cite{yi2024survey, brooks2023large}, sequential decision-making\cite{zhao2024expel, zhai2025enhancing}, and coding\cite{roziere2023code, yan2023codescope, yang2024qwen2}.
By leveraging the extensive and diverse human knowledge acquired during the pre-training phase, LLMs possess a broad range of world knowledge and emergent reasoning abilities\cite{yao2023react, guo2025deepseek}.
Consequently, LLM-Augmented RL has emerged as a promising approach to overcome the challenges faced by traditional RL methods\cite{cao2024surveylargelanguagemodelenhanced}.
To tackle the issues of long-term planning and sparse rewards, previous LLM-based approaches have evolved into two main paradigms: LLM as a planner\cite{brohan2023can, dalal2024plan, du2023guiding} and LLM as a rewarder\cite{xie2024text2reward, ma2023eureka, hazra2024revolve}, as illustrated in Figure \ref{fig:overview}.
As a planner, LLM decomposes complex, long-term tasks into a series of short-term, goal-oriented subtasks, incorporating subtask priors into the decision-making process.
As a rewarder, LLM constructs a dense reward function for sparse-reward tasks, providing guidance to the underlying policy based on the meaningful priors encoded within the LLM.

\begin{figure}[h]
  \centering
  \includegraphics[width=\textwidth]{figures/overview.png}
  \caption{Quick view. (a), (b) and (c) shows the structual similarities and differences between LLM-Augmented methods. (d) is the average evaluation results on Metaworld\cite{yu2019meta} and Maniskill\cite{taomaniskill3}.}
  \label{fig:overview}
\end{figure}

However, LLM as a planner primarily focuses on high-level planning.
In the absence of LLM-augmented dense reward guidance, task decomposition alone—relying solely on natural-language task priors—is insufficient to support the effective learning of subtasks\cite{prasad2023adapt, dalal2024plan}.
Previous approaches have alleviated this limitation by constructing pre-trained policy libraries\cite{brohan2023can} or incorporating traditional motion planning to link subtasks\cite{dalal2024plan}, yet these methods lack flexibility and efficiency due to their non-end-to-end nature.
On the other hand, LLM as a rewarder provides low-level guidance.
However, without subtask decomposition, constructing a reward function for an entire long-term, complex task remains a significant challenge for LLMs\cite{sarukkai2024automatedrewardsllmgeneratedprogress, cao2024drlc}.
Previous methods have enhanced the quality of LLM-generated reward functions by incorporating human feedback\cite{xie2024text2reward} or evolutionary search algorithms\cite{ma2023eureka, hazra2024revolve}, but these approaches still require extensive search and significant human expert assistance.
As shown in Figures \ref{fig:overview}a and \ref{fig:overview}b, both of these approaches focus on a single aspect of RL, neglecting the importance of integrating both high-level planning and low-level reward shaping.

To address the challenges of long-term planning and sparse rewards in reinforcement learning more efficiently, we propose a novel framework, Progress Reward Model for Reinforcement Learning (PRM4RL), which integrates both high-level long-term planning and low-level dense rewards to provide a comprehensive solution, as illustrated in Figure \ref{fig:overview}c.
In high-level planning, the task is decomposed into subtasks, each with termination conditions and fine-grained progress evaluations.
For low-level reward generation, we utilize the progress to construct an efficient dense reward signal with theoretically grounded optimality and convergence guarantees.
In practice, our pipeline decomposes a complex task into a series of manageable, goal-oriented subtasks, each associated with a natural language subtask prior.
For each subtask, we build a fine-grained progress function that dynamically tracks the execution of the task.
Inspired by the potential field theorem, we regard the progress function as a potential function and design a Progress Reward Model with optimality and convergence guarantees.
The evaluation results in Metaworld\cite{yu2019meta} and Maniskill\cite{taomaniskill3} as shown in Figure \ref{fig:overview}d demonstrate the effectiveness of our proposed framework, with significant improvements in success rate and reduced LLM calls, outperforming existing methods in both performance and efficiency.
